\documentclass[10pt,a4paper]{amsart}
\usepackage[margin=19mm]{geometry}
\usepackage{amsmath,amssymb,mathtools}
\usepackage[scr=rsfs,cal=euler]{mathalfa}
\usepackage{xfrac}
\usepackage[unicode,hidelinks,bookmarks=false]{hyperref}
\newcommand{\ie}{i.e.\ }
\newcommand{\cf}{cf.\ }
\newcommand{\resp}{respectively\ }
\newcommand{\Subsection}{\subsection}
\providecommand{\nobreakdash}{\nobreak\hbox{-}}
\setlength{\emergencystretch}{2em}
\multlinegap0pt
\usepackage{bm}
\usepackage{bbm}
\usepackage{dsfont}
\usepackage{xspace}
\usepackage{cancel}
\usepackage{mathtools}
\usepackage{amsthm}
\makeatletter
\@ifundefined{thm}{\newtheorem{thm}{Thm}}{}
\@ifundefined{lem}{\newtheorem{lem}[thm]{Lemma}}{}
\makeatother
\def\fr{\mbox{$\frac{1}{2}$}}
\def\rd{{\rm d}} \def\re{{\rm e}} \def\fr{\mbox{$\frac{1}{2}$}}
\def\uh{\widehat {u}} \def\Zh{\widehat {Z}}
\title[Heteroclinic connections between finite-amplitude periodic o]{Heteroclinic connections between finite-amplitude periodic orbits emerging from a codimension two singularity}
\author{Thomas J. Bridges, David J.B. Lloyd, Daniel J. Ratliff, Patrick Sprenger}
\date{Source: arXiv:2602.05680v1 (CC BY 4.0)}
\begin{document}
\maketitle

\noindent\emph{Source and licence.} Heteroclinic connections between finite-amplitude periodic orbits emerging from a codimension two singularity. Version arXiv:2602.05680v1 (CC BY 4.0), distributed under the Creative Commons Attribution 4.0 International licence, as recorded in the arXiv licence metadata for this exact version and shown on the abstract page https://arxiv.org/abs/2602.05680v1. Full rights notice: licenses/arxiv-2602.05680-CC-BY-4.0.txt.\par\smallskip

\noindent\textit{Editorial scope.} Counted unit: the Lem whose source label is \texttt{lemma-Hk} in the article (source0.tex lines 1007--1011, source label \texttt{lemma-Hk}), delivered together with the article's own complete original proof of it (source0.tex lines 1013--1052, one \texttt{proof} environment of 40 lines). No mathematics is written, repaired or re-derived: statement and proof are the article's own text line for line. Required context retained uncounted: primary-ode (lines 218--220); H-def (lines 234--236); A-per (lines 952--954). Every definition, display and label of those lines is retained unchanged. Omitted from this package and not redistributed: 1--217, 221--233, 237--951, 955--1006, 1012--1012, 1053--2427. Nothing inside a retained excerpt is omitted or altered: every retained body is delivered exactly as the article's lines are written, so no omission receipt of any kind accompanies this package. Citations: the retained excerpts carry no citation command at all, so no bibliography entry of the article is transcribed and no reference is redistributed. Reference adaptation: none was needed and none was made; every reference inside the delivered unit resolves inside it, so no unresolved reference or citation is produced. Printed slips kept literal: no printed word, symbol, hypothesis or number of the counted statement or of its proof was altered. Layout and class: the article's own class is \texttt{article} with packages including bbm, bbold, amsmath, amsthm, amssymb, mathrsfs, graphicx, authblk, framed, hyperref, geometry, comment, xcolor, inputenc; the wrapper is the lane's pinned \texttt{amsart} editorial wrapper at 10pt on A4, which restates the theorem-like environments the retained text needs and adds the article's own macro definitions that the retained text needs (\texttt{\textbackslash fr}, \texttt{\textbackslash rd}, \texttt{\textbackslash uh}), with extra wrapper packages (bm, bbm, dsfont, xspace, cancel, mathtools). The delivered numbering is the unit's own. Assets: no retained span contains a figure environment, a graphics-inclusion command, a PDF-page inclusion, a table or a TikZ picture; no image file is copied into this package, the package declares no asset dependency and no image file is redistributed with it. Licence: version arXiv:2602.05680v1 of this article is distributed under the Creative Commons Attribution 4.0 International licence, as recorded in the arXiv licence metadata for this exact version and shown on the abstract page https://arxiv.org/abs/2602.05680v1. A full rights notice is delivered at licenses/arxiv-2602.05680-CC-BY-4.0.txt. Characters: every retained excerpt is pure ASCII, so the delivered engine, XeLaTeX with its default Latin Modern text font, reports no missing character.
\begin{equation}\label{primary-ode}
u_{xxxx} + \sigma u_{xx} + V'(u) = 0 \,,
\end{equation}

\begin{equation}\label{H-def}
H = u_xu_{xxx} - \fr u_{xx}^2 + \fr \sigma u_x^2 + V(u)\,.
\end{equation}

\begin{equation}\label{A-per}
  A(k) = \frac{1}{2\pi}\int_0^{2\pi} \big( \sigma k \uh_z^2 - 2k^3 \uh_{zz}^2\big)\,\rd z\,,
\end{equation}

\begin{lem}\label{lemma-Hk} For the equation (\ref{primary-ode}) with Hamiltonian $H$ in (\ref{H-def}) and action (\ref{A-per}) evaluated on a family of periodic orbits with $k\neq0$,
    \begin{equation}\label{HkAk}
H_k = k A_k\,.
\end{equation}
 \end{lem}

\begin{proof}
    First, since the value of $H$ is independent of $z$, we have
%
\begin{equation}\label{Hk-per}
H = \frac{1}{2\pi}\int_0^{2\pi} H\,\rd z =
\frac{1}{2\pi}\int_0^{2\pi} \left( -\frac{3}{2}k^4\uh_{zz}^2+
\fr\sigma k^2 \uh_{z}^2 + V(\uh)\right)\,\rd z \,,
\end{equation}
%
after integration by parts. Differentiate this expression with respect to $k$,
\[
H_k =
\frac{1}{2\pi}\int_0^{2\pi} \left( - 6k^3\uh_{zz}^2
-3k^4\uh_{zz}\uh_{zzk} +\sigma k \uh_{z}^2
+\sigma k^2 \uh_{z}\uh_{zk}+ V'(\uh)\uh_k\right)\,\rd z \,.
\]
Rearrange, using the fact that $\uh$ is a solution
\[
H_k =
\frac{1}{2\pi}\int_0^{2\pi} \left(\sigma k \uh_z^2
+2\sigma k^2 \uh_z\uh_{zk} - 6k^3\uh_{zz}^2 -4k^4\uh_{zz}\uh_{zzk}
+ \Big[k^4 \uh_{zz}\uh_{zzk}- \sigma k \uh_z\uh_{zk}
+ V'(\uh)\uh_k\Big]\right)\,\rd z \,.
\]
The term inside the square brackets vanishes on solutions (after integration by
parts), leaving
\[
H_k =
\frac{1}{2\pi}\int_0^{2\pi} \left(\sigma k \uh_z^2
+2\sigma k^2 \uh_z\uh_{zk} - 6k^3\uh_{zz}^2 -4k^4\uh_{zz}\uh_{zzk}
\right)\,\rd z\,.
\]
Now differentiate $A(k)$ in (\ref{A-per}),
\[
A_k = \frac{1}{2\pi}\int_0^{2\pi}\left(
\sigma \uh_z^2 + 2\sigma k \uh_z\uh_{zk} -6k^2\uh_{zz}^2 - 4k^3\uh_{zz}\uh_{zzk}\right)\,\rd z\,.
\]
Comparing the two proves the result. That is, when $k\neq0$, we have
$kA(k)=H(k)$.
\end{proof}

\end{document}
